
\section{Photovoltaic Integration and Forecasting}

Photovoltaic (PV) systems have become the cornerstone of residential renewable generation due to their modularity, decreasing capital costs, and policy incentives. The efficiency of commercial silicon-based solar cells has seen steady improvement, now routinely exceeding 25\% \cite{Green2022}, making rooftop installations an economically viable option for homeowners. However, the integration of high penetrations of PV into the distribution grid is not without challenges.

The fundamental limitation of solar energy is its intermittency. Power generation is strictly diurnal and highly sensitive to meteorological conditions such as cloud cover, aerosols, and humidity. This variability introduces significant uncertainty into the energy management problem. An effective EMS relies on accurate knowledge of future generation to schedule battery charging and discharging cycles optimally. Consequently, the performance of the control strategy is intrinsic to the quality of the solar forecast.

Forecasting techniques generally fall into three categories: physical models, statistical models, and machine learning approaches. Physical models rely on Numerical Weather Prediction (NWP) data and mathematical models of the PV panel's electrical characteristics. While useful for long-term planning, they often lack the spatial and temporal resolution required for real-time microgrid control \cite{Theocharides2023}. Statistical methods, such as Auto-Regressive Integrated Moving Average (ARIMA), use historical time-series data to predict future values but struggle to capture non-linear weather patterns.

Recent research has focused on machine learning and deep learning models for short-term forecasting. Studies have demonstrated that techniques such as Long Short-Term Memory (LSTM) networks can achieve higher accuracy by learning complex temporal dependencies \cite{Deanseekeaw2024}. Despite these advancements, no forecast is perfect. "Sim-to-Real" gaps, defined as the discrepancy between the forecasted profile used for planning and the actual realized generation, remain unavoidable. This reality necessitates the development of robust control strategies capable of absorbing prediction errors without catastrophic performance degradation, a core motivation for the R-SAC agent proposed in this thesis.
